\documentclass[11pt,aspectratio=169]{beamer}

\usetheme{Madrid}
\usecolortheme{default}

\usepackage{amsmath,amssymb,mathtools}

\setbeamertemplate{navigation symbols}{}
\setbeamertemplate{footline}[frame number]

\title{What is the type signature of learning?}
\author{Vinayak Pathak}
\date{\today}

\begin{document}

\begin{frame}[plain,noframenumbering]
    \titlepage
\end{frame}

\begin{frame}{Learning builds agents}
\Large
\begin{itemize}[<+->]
    \item Today's agents are built using learning.
    \item Future agents probably will be too.
    \item We need learning processes that produce safe agents.
\end{itemize}
\end{frame}

\begin{frame}{What is learning?}
\Large
\begin{itemize}[<+->]
    \item The Bayesian answer:
    \[
        (\text{Belief},\text{Evidence})\longmapsto\text{Belief}
    \]
    \item Objection 1: Can we compute the update?
    \item Objection 2: What does the update guarantee?
\end{itemize}
\end{frame}

\begin{frame}{Objection 1: computation}
\large
\begin{itemize}
    \item<1-> Domain: $x\in\{0,1\}^n$.
    \item<2-> Query access to $p(x)$ and $p(y\mid x)$.
\end{itemize}
\uncover<3->{
\begin{block}{Toy example}
\[
    p(x\mid y)=\frac{p(x)p(y\mid x)}{\sum_z p(z)p(y\mid z)}
\]
\end{block}
}
\uncover<4->{
\begin{alertblock}{Question}
Can we compute it exactly with $\operatorname{poly}(n)$ queries?
\end{alertblock}
}
\uncover<5->{
\begin{center}
\Large\alert{No.} Calculating the denominator requires $2^n$ queries.
\end{center}
}
\end{frame}

\begin{frame}{Objection 2: guarantees}
\large
\begin{itemize}
    \item<1-> Cox's axioms motivate coherent belief updates.
    \item<2-> They do not choose the performance guarantee.
\end{itemize}
\begin{columns}[T,onlytextwidth]
\begin{column}{0.48\textwidth}
\uncover<3->{
\begin{block}{Bayes}
\[
    \min_a\;\mathbb{E}_{\theta\sim p} L(a,\theta)
\]
Minimize prior-averaged loss.
\end{block}
}
\end{column}
\begin{column}{0.48\textwidth}
\uncover<4->{
\begin{alertblock}{Worst case}
\[
    \min_a\;\max_\theta L(a,\theta)
\]
Minimize the largest loss.
\end{alertblock}
}
\end{column}
\end{columns}
\end{frame}

\begin{frame}{A small average can hide a large failure}
\Large
\begin{itemize}
    \item<1-> Easy case: probability $1-\varepsilon$, loss $0$.
    \item<2-> Hard case: probability $\varepsilon$, loss $1$.
\end{itemize}
\begin{columns}[T,onlytextwidth]
\begin{column}{0.48\textwidth}
\uncover<3->{
\begin{block}{Average loss}
\centering
\Huge $\varepsilon$
\end{block}
}
\end{column}
\begin{column}{0.48\textwidth}
\uncover<4->{
\begin{alertblock}{Worst-case loss}
\centering
\Huge $1$
\end{alertblock}
}
\end{column}
\end{columns}
\end{frame}

\begin{frame}{When can Bayes catch up?}
\Large
\begin{itemize}
    \item<1-> The truth lies in the prior's support.
    \item<2-> The evidence identifies the truth.
    \item<3-> Updating continues long enough.
\end{itemize}
\uncover<4->{
\[
    p(\text{truth}\mid\text{evidence})\longrightarrow 1
\]
}
\uncover<5->{
\begin{center}
Only then can average-case behavior approach the optimum for the true world.
\end{center}
}
\end{frame}

\begin{frame}{Finite computation blocks the escape hatch}
\Large
\begin{itemize}[<+->]
    \item Exact updating can be intractable.
    \item A tractable learner has finite time and may need approximations.
    \item We cannot rely on infinite posterior concentration.
    \item Small Bayesian risk does not imply small worst-case loss.
\end{itemize}
\end{frame}

\end{document}
